% Analytic illustration only: identified interval [1,3], outer radius 2/sqrt(n).
\begin{tikzpicture}
  \begin{axis}[
    width=\linewidth,height=61mm,
    xmode=log,log basis x=2,xmin=4,xmax=256,
    ymin=0,ymax=4.35,xtick={4,16,64,256},xticklabels={4,16,64,256},
    ytick={0,1,2,3,4},axis lines=left,
    xlabel={样本量$n$},ylabel={价值（回报）},
    tick label style={font=\figurefont\fontsize{8}{10}\selectfont},
    label style={font=\figurefont\fontsize{8}{10}\selectfont},
    axis line style={BookMuted,line width=.7pt},
    tick style={BookMuted},clip=false
  ]
    \path[fill=BookTeal!8] (axis cs:4,1) rectangle (axis cs:256,3);
    \addplot[BookTeal,dashed,line width=.9pt,domain=4:256] {1};
    \addplot[BookTeal,dashed,line width=.9pt,domain=4:256] {3};
    \addplot[BookBlue,line width=1.1pt,domain=4:256,samples=80]
      {1-2/sqrt(x)};
    \addplot[BookBlue,line width=1.1pt,domain=4:256,samples=80]
      {3+2/sqrt(x)};
    \node[font=\figurefont\fontsize{8}{10}\selectfont,text=BookInk]
      at (axis cs:30,2) {不可消除的识别宽度};
    \draw[{Latex[length=1.6mm]}-{Latex[length=1.6mm]},
      BookAmber,line width=.8pt]
      (axis cs:8,3.04) -- (axis cs:8,3.66);
    \node[anchor=west,font=\figurefont\fontsize{7.5}{9}\selectfont,
      text=BookAmber] at (axis cs:9,4.02) {统计误差示意};
  \end{axis}
\end{tikzpicture}
